% Partie : méthodes
% Chapitre : démonstrations
% Section : introduction

\documentclass[../../../include/open-logic-section]{subfiles}

\begin{document}

\olfileid[fr]{mth}{prf}{int}

\olsection{Introduction}

Votre expérience de la logique élémentaire
vous a peut-être familiarisé avec un système
de !!{derivation}, probablement en
déduction naturelle, en notation de Fitch
ou sous une autre forme, ou peut-être
avec des arbres de preuve. Vous vous souvenez
sans doute avoir construit des démonstrations
dans ces systèmes, soit pour établir
!!a{formula}, soit pour montrer la validité
d'un argument. Pour y parvenir, vous
appliquiez les règles du système jusqu'à
obtenir le résultat voulu. Lorsque nous
raisonnons \emph{sur} la logique, nous
démontrons aussi des énoncés, mais,
le plus souvent, sans utiliser de
système de !!{derivation}. En effet,
la plupart des démonstrations que nous
étudierons sont rédigées en langue
ordinaire, éventuellement avec quelques
symboles, plutôt qu'entièrement dans
le langage de la logique du premier ordre.
Au début, construire une telle démonstration
peut vous laisser perplexe : comment
démontrer quelque chose sans système
de !!{derivation} ? Par où commencer ?
Comment savoir si votre démonstration
est correcte ?

Avant de tenter une démonstration, il
importe de savoir ce qu'elle est et
comment la construire. Comme son nom
l'indique, une \emph{démonstration}
sert à établir la vérité d'une
affirmation. Imaginez un dialogue :
quelqu'un vous demande, par exemple,
si tout nombre premier autre que
deux est impair. Répondre «~oui~»
ne suffit pas ; cette personne
pourrait vous demander \emph{pourquoi}.
Vous lui donneriez alors une
démonstration.

Dans la conversation courante,
une réponse esquissée ou incomplète
peut suffire. En logique et en
mathématiques, nous voulons toutefois
une démonstration rigoureuse :
nous voulons établir une affirmation
sans laisser place au \emph{moindre}
doute. Chaque étape doit donc être
justifiée, et sa justification
doit être convaincante. Par exemple,
une hypothèse invoquée doit bien
figurer parmi celles de l'énoncé
à démontrer ; les définitions
doivent être appliquées correctement ;
les inférences utilisées doivent
être valides.

En général, nous démontrons un
énoncé. Les énoncés à démontrer
portent divers noms : propositions,
théorèmes, lemmes ou corollaires.
Une proposition est un énoncé
qui mérite d'être démontré et
consigné, sans être nécessairement
profond ni souvent utilisé.
Un théorème est une proposition
importante. Sa démonstration
se décompose souvent en plusieurs
étapes. Il porte parfois le nom
de la personne qui l'a démontré
la première, comme le théorème
de Cantor ou celui de
L\"owenheim--Skolem. Il peut
aussi être nommé d'après le
fait qu'il établit, comme
le théorème de complétude.
Un lemme est une proposition
ou un théorème utilisé pour
démontrer un résultat plus
important. Il arrive cependant
que des lemmes soient importants
en eux-mêmes et portent aussi
le nom de leur auteur, comme
le lemme de Zorn. Un corollaire
est un résultat qui découle
facilement d'un autre.

Un énoncé à démontrer comporte
souvent des suppositions qui
précisent les objets dont il
est question. Il peut commencer
par «~Soit $!A$ !!a{formula}
de la forme $!B \lif !C$~»
ou «~Supposons
$\Gamma \Proves !A$~»,
par exemple. Ce sont les
\emph{hypothèses} de la
proposition, du théorème
ou du lemme ; vous pouvez
les tenir pour vraies dans
votre démonstration. Elles
délimitent ce que nous
démontrons et nomment
certains objets du raisonnement.
Ainsi, si votre proposition
commence par «~Soit $!A$
!!a{formula} de la forme
$!B \lif !C$~», elle ne
porte que sur les formules
d'un certain type, à savoir
les implications. On comprend
que $!B \lif !C$ désigne
une implication quelconque
sur laquelle portera la
démonstration.

\end{document}
